\chapter{Arquitetura em Quatro Camadas}

\section{Estrutura de Abstração}
A orquestração Quantiflux divide o fluxo de trabalho em quatro níveis distintos e complementares:

\subsection{Nível Conceitual}
Representa a tese e a modelagem formal do sistema. Define as entidades fundamentais, os contratos imutáveis e as invariantes de domínio. Não contém detalhes de implementação ou bibliotecas específicas.

\subsection{Nível Estratégico}
Converte a tese conceitual no Plano Diretor de Engenharia (\emph{Master Engineering Plan}). Responsável por decompor o projeto em fases lógicas e ordenar as dependências entre módulos.

\subsection{Nível Tático}
Atua como o gerenciador de contexto e despachante. Filtra as informações relevantes do Plano Diretor para a Onda corrente e dispara subagentes especialistas via primitivas como \texttt{delegate\_task}.

\subsection{Nível Operacional}
Constituído pelos operadores e executores de código. Trabalham em ambientes limpos, escrevendo testes unitários e implementações locais. Submetem seus artefatos para validação nos portões de qualidade.